\subsection{The ICM programme: the precise historical comparison}
\begin{proposition}[Known, with the primary formulations compared]\label{cat:J9}
The ``mysterious functor'' lead is correct for reconstruction of
Barsotti--Tate isogeny data through crystalline representations.
It does not identify Fontaine's crystalline category with the literal
all-Hodge--Tate category of 161-5/023.
The relevant ICM article is \emph{Groupes de Barsotti-Tate et cristaux},
Nice 1970, tome~1, pp.~431--436 (published in 1971), not the separately
published Montreal course whose title adds \emph{de Dieudonn\'e}.
\end{proposition}
\begin{proof}[Numbered comparison of the opened statements]
\begin{enumerate}
\item Grothendieck's \S5, p.~435, Remarks~2--3
\cite{catGrothendieckICM1970}, ask for the Galois representation from
the filtered Dieudonne space and for the converse reconstruction and
essential image.  Fontaine's crystalline equivalence in
\cite[\S5.2]{catFontaine1982}, together with its BT compatibility
in \S6.2, provides these inverse constructions on the admissible
category.  One must choose consistently between covariant and
contravariant Dieudonne modules; the reconstruction assertion is
unchanged by the corresponding dualization.
\item The intrinsic description of that essential image requires a
further theorem.  Colmez--Fontaine's Theorem~A and \S4.1 identify
weak admissibility with admissibility
\cite{catColmezFontaine2000}.  Kisin's Corollary~(2.2.6)
\cite{catKisin2006} identifies the crystalline weight-$[0,1]$
representations with rational Tate modules of BT groups.  These
statements together answer the BT essential-image problem at the
isogeny level for the perfect-residue-field scope stated above.
Theorem~\ref{cat:J7} then supplies the precise uniqueness of the
tensor extension from actual BT objects with the specified graded
comparison.  It is not a classification by Newton slopes alone.
\item This connection is also documented by the authors of the modern
theorem: Colmez--Fontaine's Introduction, p.~3 of the opened author
version, identifies Grothendieck's questions on the ``foncteur
myst\'erieux'' as an origin of the theory, and its reference~[23] is
exactly the ICM article above.  Thus the historical link does not
depend on the later study's assertion alone.
\item The final paragraph of ICM p.~435 asks for analogous
comparisons in higher cohomology and over more general complete bases.
Those broader requests are not all the assertion of the 1982 BT
comparison theorem.  Likewise, the S3 obstruction in
Theorem~\ref{cat:J8} proves that the larger category literally written
in the dossier cannot simply be substituted into the modern theorem.
The defensible identification is of the reconstruction programme and
its crystalline solution, with these categorical hypotheses retained.
\end{enumerate}
\end{proof}

